% ========== 2.4 FUNCTIONAL REQUIREMENTS ==========
% Core system capabilities and features
\clearpage
\section{2.4 Functional Requirements}
\addcontentsline{toc}{section}{2.4 Functional Requirements}
\label{sec:functional-requirements}

\lettrine{F}{unctional requirements} specify the capabilities that Symphony must provide to support user needs and achieve project objectives. These requirements define observable system behaviors and interactions from a user perspective, directly corresponding to the system features presented in Section 2.2. Each functional requirement category addresses specific capabilities necessary for implementing Symphony's core features.

% ========== FR Category 1: AI Orchestration & Intelligence ==========
\subsection*{2.4.1 AI Orchestration \& Intelligence}
\addcontentsline{toc}{subsection}{2.4.1 AI Orchestration \& Intelligence}
\label{sec:fr-ai-orchestration}

\textbf{FR-1.1 Intelligent Agent Coordination:}
The Conductor shall coordinate multiple specialized AI agents autonomously, managing task decomposition, agent selection, execution sequencing, and inter-agent communication without manual intervention.

\textbf{FR-1.2 Natural Language Understanding:}
The system shall accept natural language prompts describing development objectives and translate them into executable workflows with appropriate agent orchestration.

\textbf{FR-1.3 Context-Aware Processing:}
The system shall maintain project context, existing codebase patterns, coding standards, and user preferences to inform orchestration decisions and ensure consistency.

\textbf{FR-1.4 Task Decomposition:}
The Conductor shall automatically decompose complex development tasks into manageable subtasks, assigning each to the most appropriate specialized agent.

\textbf{FR-1.5 Parallel Execution Optimization:}
The system shall identify opportunities for parallel agent execution, optimizing workflow performance while respecting task dependencies.

\textbf{FR-1.6 Dependency Resolution:}
The DAG Tracker shall automatically resolve workflow dependencies, determine optimal execution order, and manage inter-task relationships.

% ========== FR Category 2: Visual Workflow Management ==========
\subsection*{2.4.2 Visual Workflow Management}
\addcontentsline{toc}{subsection}{2.4.2 Visual Workflow Management}
\label{sec:fr-workflow}

\textbf{FR-2.1 Node-Based Workflow Composition:}
The Harmony Board shall provide intuitive node-based visual interface for creating custom workflows by connecting AI agents and operations without programming knowledge.

\textbf{FR-2.2 Workflow Execution:}
The system shall execute defined workflows with real-time monitoring, progress tracking, intermediate result visualization, and execution state management.

\textbf{FR-2.3 Workflow Templates (Melodies):}
The system shall support saving, loading, and sharing workflow templates through the marketplace, enabling reuse and knowledge transfer across projects and teams.

\textbf{FR-2.4 Configurable Parameters:}
Workflows shall support configurable parameters for fine-tuning agent behavior, setting resource limits, defining timeouts, and specifying error handling strategies.

\textbf{FR-2.5 Real-Time Monitoring:}
The Harmony Board shall provide real-time visualization of workflow execution, agent activity, and system performance during workflow runs.

% ========== FR Category 3: Extension System ==========
\subsection*{2.4.3 Extension System}
\addcontentsline{toc}{subsection}{2.4.3 Extension System}
\label{sec:fr-extensions}

\textbf{FR-3.1 Extension Categories:}
The system shall support three extension categories: Instruments (AI Models), Operators (Workflow Utilities), and Addons (Interface Enhancements), each with appropriate isolation and integration mechanisms.

\textbf{FR-3.2 Hot-Swappable AI Models:}
The system shall support integration of multiple AI models (GPT, Claude, custom models) as interchangeable Instrument extensions, enabling model selection based on task requirements.

\textbf{FR-3.3 Extension Installation and Management:}
The system shall support extension discovery, installation, configuration, updates, and removal through integrated marketplace interface with automatic dependency resolution.

\textbf{FR-3.4 Extension Isolation:}
The system shall execute user-facing extensions in isolated processes with sandboxed environments, preventing malicious code from affecting core system stability.

\textbf{FR-3.5 Extension Lifecycle Management:}
The system shall manage extension activation, deactivation, and hot-reloading without requiring system restart, maintaining system availability during updates.

\textbf{FR-3.6 Extension Marketplace:}
The system shall host marketplace for extension distribution supporting multiple business models (free, freemium, pay-per-use, enterprise licensing) with revenue sharing mechanisms.


% ========== FR Category 4: Error Handling & Reliability ==========
\subsection*{2.4.4 Error Handling \& Reliability}
\addcontentsline{toc}{subsection}{2.4.4 Error Handling \& Reliability}
\label{sec:fr-reliability}

\textbf{FR-4.1 Fault Isolation:}
The microkernel architecture shall isolate extension failures, preventing individual component crashes from causing system-wide failures or data loss.

\textbf{FR-4.2 Automatic Retry Mechanism:}
The system shall implement automatic retry with exponential backoff for transient failures, configurable retry limits, and intelligent failure classification.

\textbf{FR-4.3 Alternative Agent Selection:}
Upon agent failure, the Conductor shall automatically select alternative agents capable of completing the task, maintaining workflow progress without user intervention.


\textbf{FR-4.4 Failure Learning:}
The Conductor shall learn from failure patterns, adjusting orchestration strategies to prevent recurrence and improving system reliability through experience.

\textbf{FR-4.5 Health Monitoring:}
The system shall continuously monitor component health, detect anomalies, and trigger proactive recovery actions before failures impact user workflows.

